\section{Discussion}\label{sec:discussion}

\paragraph{Deployment guidance.} For reward-model verifiers of 4B parameters and above,
4-bit weight quantization (GPTQ, AWQ or NF4) is a safe default for Best-of-$N$
selection on mathematical problems: across three preregistered settings, two
generators, an ORM family and a PRM, the accuracy change stayed within about 1.5\pp\
and the verifier's advantage over majority vote was largely preserved. Under a fixed
memory budget, the freed memory is better spent on a larger quantized verifier than on
a smaller full-precision one (c2). Three bits is not a safe default: it is harmless in
some regimes (H6) but costs 2--4\pp\ in others (H7, main study) and is catastrophic for
the 1.7B ORM, and the smallest verifier is also sensitive to the quantization method
(RTN versus GPTQ).

\paragraph{Why quantization rarely matters, and when it could.} Best-of-$N$ only needs
the top of the ranking to keep a correct answer, and candidate pools are redundant:
many candidates share the same final answer. Quantization perturbs scores enough to
change the selected candidate on a quarter to a third of problems, yet the new winner
usually carries the same answer. The risk is concentrated on contested problems, where
competing answers have similar support and the verifier's preference is what decides;
there, coarse quantization does tip choices (3 bits), and in our H6 and H7 data the
perturbation tends to move the selection towards the consensus answer. This suggests
that quantization damage should grow with the share of contested problems and with how
much a verifier's decisions depart from the consensus, rather than with nominal
difficulty or with $N$, which is consistent with our unsupported difficulty and budget
hypotheses (c4, c5).

\paragraph{How much a verifier is worth.} The question ``does quantization destroy the
value a verifier adds?'' presupposes that the verifier adds value. With the weak
generator the Skywork ORMs were within about one point of majority vote overall, and
with a strong math generator on hard problems majority vote matched the PRM and beat the
ORM by 3.8\pp. Only in a regime selected for verifier advantage (H7) did both verifiers
clearly beat majority vote. Practitioners should measure a verifier against
self-consistency on their own generator before paying its memory and latency cost.

\paragraph{Evaluating selection.} Three practices changed our conclusions or their
size. (i) A strict answer grader turned a reward-model preference for prose answers into
apparent quantization ``gains'' of up to 8\pp; a format/content decomposition and a
lenient grader removed them, and exhaustive blinded human grading reduced the registered
memory-frontier effect from 2.25 to 1.4\pp. (ii) A fixed-precision null control (same
weights, different batch layout) changes 5--15\% of winners across our studies; winner-change rates without
such a floor overstate instability. (iii) Development-set signals that motivated two
registered hypotheses did not replicate on 1,200 evaluation problems; preregistering
them made the non-replication visible rather than leaving it in a file drawer.

\paragraph{Limitations.} The ORMs come from a single family (Skywork-Reward-V2, built on
Qwen3), the PRM is a Qwen model and the generators are Qwen and Llama models, so
family-specific robustness cannot be excluded. Accuracy experiments use simulated
quantization; packed 4-bit kernels agreed with it on 98.8--99.2\% of selected answers,
but we measured memory, not latency or throughput. The tasks are GSM8K and MATH, whose
test problems may have been seen during pre-training of the generators and verifiers;
contamination affects BF16 and quantized verifiers equally because both score the same
frozen pools, but it can change pool redundancy and verifier advantage. Candidates were
capped at 1,024 or 2,048 tokens, and 6.5--19\% reached the cap, although excluding them
does not change any conclusion. The value-retention result is confirmatory for the PRM
and only narrowly so for the ORM, and we did not find a strong-generator regime in which
the verifiers clearly beat majority vote, so retention of a large advantage under a
strong generator remains untested. The analysis plans were fixed in version control
before the corresponding data were generated or scored; we did not use an external
registry.
